\section{Deferring separator work until the scalar query needs it}
\label{sec:adaptive-potts}

The separator backend supplies a useful universal bound, but eagerly building
its certificate slows many easy scalar queries.  The new optional
\code{potts-adaptive} backend starts with the ordinary order and watches
\emph{actual completed transitions}.  It postpones separator construction
until that work becomes large, while retaining the preceding section's
subexponential fixed-color scalar bound.  This is the introsort principle
applied to one exact invariant query: limit speculative work, then use a
method with a proved bound.  It is not yet an analogous worst-case bound
for the complete unknot recognizer.

\subsection{The policy and its persistent state}

For an $n$-crossing input, let the initial work threshold be $B=64n$ and the
short-tail allowance be $r_0=2\lceil\log_2(n+1)\rceil$.  The evaluator begins
on a supplied order, or the existing greedy order when none is supplied.
At each completed crossing before the last, it observes the cumulative
transition count.  The actions are as follows.
\begin{enumerate}
\item Below $B$, continue normally.  A completed scalar answer needs no
separator certificate.
\item On first reaching $B$, if at most $r_0$ crossings remain, finish the
current table.  This enters \code{polynomial-tail} mode; the argument below
supplies a bound without requiring a width certificate.
\item Otherwise construct and independently verify a separator hierarchy.
Compare its order with the better of the current order and a freshly
computed greedy order, using (maximum frontier, total frontier).
\item If the selected order is unchanged, retain the \emph{same} Potts table
and continue.  No processed crossing or transition is repeated.  This is
\code{certified-continue} mode.
\item A changed order enters \code{certified-restart} mode.
Discard the unfinished scalar table and restart once.  The old exception handler
is exited before allocating the replacement table, releasing the abandoned
evaluator's traceback and local references.
\end{enumerate}

A local state-cap failure can also trigger the preparation when it has not
already happened and transition allowance remains.  If the selected order
is unchanged, repeating the query cannot repair that failure: publish the
completed order certificate and decline immediately.  A spent transition
allowance never triggers preparation.  There is at most one separator
preparation and at most one restart.  A fixed state cap may still interrupt
the polynomial tail, in which case the same rules apply.

The \code{ordering\_policy} evidence records the mode, threshold, short-tail
allowance, preparation/restart counts, abandoned crossings and transitions,
and total spent transitions.  \code{PottsLimit} carries work counters only;
it never exposes a partial partition function as a completed value.  The
public exact evaluator's existing complete-result behavior is unchanged.

\subsection{A single transition budget, including discarded work}

Suppose the user supplies transition allowance $T$, and the first attempt
spends $a$ transitions.  A restarted evaluator receives exactly $T-a$.
If it spends another $b$, the returned count and any exhaustion exception
report $a+b\le T$.  There is no fresh allowance for the second implementation.
State caps keep their existing meaning as limits on represented table keys;
the reported peak is the larger attempt peak, not a byte-memory estimate.
Other scalar geometry and coefficient statistics describe the final attempt.

Order preparation is charged to the shared global deadline, not to the
scalar-transition counter.  The deadline callback propagates through the
policy and verifier.  A preparation interrupted globally publishes no partial
certificate.  After a complete witness is available, a later local scalar
failure retains it for the recognition fallback.  Subsequent crossing-changing
Reidemeister simplification rebuilds it on the reduced diagram, as in the
preceding section.  When the scalar query finishes in ordinary or
polynomial-tail mode, no separator witness has been computed; those modes
make no width promise for the later homology scan.

\subsection{Why the deferral preserves the asymptotic scalar bound}

Fix the number of colors $q\ge5$.  One crossing introduces at most two Tait
vertices.  From one stored spin-equality state there are at most $q^2$
extensions, because every canonical extension represents at least one
labeled choice for those two vertices.  Aggregation and cancellation only
decrease the number of states.  Moreover, after at least one crossing the
number of stored states is at most the number of transitions spent so far.

Immediately before the first completed stage reaching $B$, at most
$B+1$ states are available.  The crossing that first passes the threshold
therefore takes at most $(B+1)q^2$ further transitions.  Stopping only at a
crossing boundary causes polynomial, rather than unbounded, overshoot.
The same observation bounds a final crossing that completes the entire query
before the policy is invoked.  Canonical-key work and exact coefficient
arithmetic have polynomial bit cost at fixed $q$.

If the policy enters a short tail of $r\le r_0$ crossings, the number of
additional transitions is bounded by
\[
 O\!\left((B+1)q^2\sum_{j=1}^{r}q^{2j}\right)
 =O\!\left(nq^{2r+O(1)}\right)
 =n^{O(1)}\quad\text{for fixed }q.
\]
The degree of this polynomial can be large.  This is an asymptotic safeguard,
not a promise that every permitted tail is practically cheap.  Existing
state, transition and deadline limits remain active.  If a state-cap failure
interrupts the tail and requests a new order, the abandoned tail work is
still bounded by this polynomial.

In every other uncapped case the policy eventually obtains a certified
$O(\sqrt n)$-frontier order.  Continuing an unchanged order or restarting
on the selected order costs $\operatorname{poly}(n)2^{O(\sqrt n)}$ bit
operations.  A continuing run does not repeat its earlier work; a restarted
run adds only the polynomially bounded first attempt.  Construction of the
hierarchy is polynomial.  Hence the complete adaptive \emph{scalar} query
has the same general subexponential upper bound as eager separator evaluation.

This theorem does not assert competitiveness with the fastest order on each
input.  Under a finite allowance, discarded transitions may prevent completion
where some single-order query would have completed.  Frontier improvement also
does not force fewer spin states after exact cancellations.  These are reasons
to expose the option and its policy evidence while retaining the existing
default, not reasons to reset the budget silently.

\subsection{Independent checks and measured scope}

The six new test cases cover actual measured-work restarts, state-cap
restarts, unchanged-order continuation, a completed short tail, exact shared
budget boundaries, globally interrupted preparation, zero budgets, invalid
options, CLI routing and the complete recognition fallback.  Twenty small
diagrams, three color counts and both shadings give 120 comparisons with an
independent whole-cube Laurent-Jones oracle.  That oracle uses neither Tait
graphs nor production ring arithmetic.  Its cases exercise policy switches
as well as ordinary completion.  The five-color blind weaving example still
falls through to Khovanov recognition rather than receiving an unknot verdict.

The benchmark uses seven shuffled rounds after an excluded warm-up, fresh
diagram construction and all order preparation inside every timer.  It
separates 15 scalar scopes from 11 full-recognition scopes.  Ordinary exact
Potts has an identical A/A control; eager separator, final adaptive and the
initial adaptive policy are compared on the same scalar queries.  Every
completed scalar comparison checks cross-multiplied normalized values,
because changing shading changes both the partition function and its
normalizing factor.  Limits remain censored outcomes.

The initial policy lacked the short-tail rule.  Its source snapshot matches
the hash in the retained initial measurements, and the final driver also
reruns it as a paired control.  On the $10\times10$ descending grid, the
threshold is crossed after crossing 90, at 6555 transitions; the whole query
requires 7109.  With ten crossings left and $r_0=14$, the final policy avoids
the otherwise unnecessary certificate.  Its median time is 21.457\,ms versus
30.556\,ms for the initial policy, a median paired improvement of 1.415 times.
On a shuffled 16-crossing grid, however, the tail rule forgoes a useful restart
and is slower than the initial policy (paired ratio 0.913).  Both results are
retained; the rule is not claimed to optimize every instance.

\begin{center}\small
\begin{tabular}{@{}lrrrrr@{}}
\toprule
Scalar input & Ordinary & A/A & Eager & Initial & Adaptive \\
\midrule
Conway & 0.499 & 0.514 & 1.134 & 0.534 & 0.547 \\
Tree medial 254 & 9.617 & 9.636 & 34.412 & 9.705 & 9.703 \\
Grid 16, shuffled & 5.875 & 5.670 & 1.390 & 5.236 & 5.692 \\
Grid 64, shuffled & limit & limit & limit & 23.877 & 23.893 \\
Grid 100 & 21.241 & 21.129 & 29.452 & 30.556 & 21.457 \\
Grid 144 & 135.286 & 135.237 & 146.759 & 149.998 & 153.190 \\
\bottomrule
\end{tabular}
\end{center}
Median milliseconds including all setup. A/A repeats ordinary exact Potts. A limit is a censored query, not a completed scalar value.


On the 254-crossing tree-medial diagram, ordinary and adaptive scalar times
are 9.617 and 9.703\,ms, versus 34.412\,ms for eager separator preparation.
On the shuffled 36-, 64- and 100-crossing grids, both adaptive policies finish
within the default budgets while ordinary and eager supplied-order queries
hit their state caps.  The adaptive fallback's fresh greedy comparison is
part of this benefit; the eager backend compares only the supplied order
with its separator order.  These capped runs have no completion-time speedup
ratio.  On the ordinary 144-crossing grid, the final policy continues the
same table after certification and takes 153.190\,ms versus 135.286\,ms for
ordinary exact Potts: the guarantee still has a preparation cost.

Full recognition on the measured inputs gains no consistent advantage over
ordinary exact Potts.  Earlier certificates bypass the scalar query on many
cases.  Avoiding eager preparation helps where it would have been used, but
this is not a general improvement over the historical matching default.
The new backend remains optional.

\subsection{A relevant geometric result, with a different scope}

The research refresh found no new delivered unknot report in the drop zone.
Lackenby's revised \emph{Some fast algorithms for curves in surfaces}
\cite{lackenby-curves2026} supplies a useful related primitive: Theorem 1.1
computes geometric intersection numbers of normal 1-manifolds in time
polynomial in triangulation size and the logarithms of their weights.
Theorem 1.2 gives a corresponding isotopy test.  The normalization theorem
also depends on the input's simplification number; that extra parameter
must not be dropped from a complexity claim.
These results strengthen the case that very large geometric objects can
admit efficient compressed operations.  They neither bound the depth of an
accelerated 3-manifold hierarchy nor the number of surviving chain objects
in our implementation.  We therefore record them as candidate geometric
primitives, without treating them as a missing recognition theorem or adding
an unused surface API to the scalar scanner.
